\chapter{先说人话：CDC 在干什么}
\label{ch:cdc-story}

\section{切面包，不是切尺子}

备份要去重：同样内容 ideally 切成同样的块。\\
固定 256KB 一切（尺子）——文件头插一个字节，后面整段边界全错位。\\
\textbf{内容定义分块}（CDC）：边界由``内容特征''决定，像在面团上找天然裂缝。

\begin{figure}[htbp]
\centering
\begin{tikzpicture}[font=\small]
  % ruler
  \node[anchor=west] at (0,2.2) {\textbf{尺子切}（固定偏移）};
  \draw[thick] (0,1.6) rectangle (11,2.0);
  \foreach \x in {0,2.75,5.5,8.25,11} {
    \draw[CutRed, very thick] (\x,1.55)--(\x,2.05);
  }
  \draw[dashed, gray] (0.4,1.6)--(0.4,0.6);
  \node[font=\footnotesize, gray] at (5.5,1.2) {插入 1 字节后：后面所有刀口全部漂移};
  % cdc
  \node[anchor=west] at (0,0.2) {\textbf{CDC 切}（内容裂缝）};
  \draw[thick] (0,-0.4) rectangle (11,0.0);
  \foreach \x in {1.8,4.2,7.0,9.6} {
    \draw[OkGreen, very thick] (\x,-0.45)--(\x,0.05);
  }
  \draw[dashed, gray] (0.4,-0.4)--(0.4,-1.2);
  \node[font=\footnotesize, gray] at (5.5,-0.85) {插入后：裂缝仍贴着相同内容图案};
\end{tikzpicture}
\caption{为什么 CDC 在增量备份里重要：刀口跟着内容走。}
\label{fig:ruler-vs-cdc}
\end{figure}

\section{三个旋钮：太碎、太大、刚刚好}

每种 CDC 都在说同一句话：``平均多久该出现一次刀口？''

\begin{metaphor}[公交车发车]
\texttt{avg\_size} 像期望发车间隔；\texttt{min} 是``刚到站不许立刻再发''；\texttt{max} 是``再不等就强制发车''。\\
算法差异在于：用什么信号决定``可以发车了''。
\end{metaphor}

\begin{figure}[htbp]
\centering
\begin{tikzpicture}[x=0.9cm, font=\small]
  \draw[thick] (0,0)--(12,0);
  \foreach \x/\t in {0/起点, 2/min, 6/期望平均, 10/max, 12/文件尾} {
    \draw (\x,0.12)--(\x,-0.12);
    \node[above, font=\footnotesize] at (\x,0.15) {\t};
  }
  \fill[blue!12] (2,-0.35) rectangle (10,0.35);
  \node[font=\footnotesize] at (6,-0.7) {只在蓝色区间里找刀口；找不到就在 max 硬切};
\end{tikzpicture}
\caption{所有后三代共用的外层``扫描窗''。}
\label{fig:scan-window-illust}
\end{figure}

\section{后三代各自的``发车信号''}

\begin{center}
\begin{tabular}{@{}llp{6.5cm}@{}}
\toprule
代 & 绰号 & 一句话 \\
\midrule
Hom & 珠串变色 & 指纹铃响了，\textbf{且}窗里珠串连通性变了，才切 \\
Chain & 密码锁 & 密码转到对齐，\textbf{且}骨架变了（可加``有环''），才切 \\
TopoPH & 双保险滤网 & 仍听指纹铃，再用三角形翻转 / 连通变化过滤 \\
Native & 星图事件 & \textbf{不听指纹铃}；内容定相``星点''上发生翻转 / $H_0$ 事件才切 \\
\bottomrule
\end{tabular}
\end{center}

\begin{pitfall}[仓库不能混切]
不同算法刀口集合不同，同一文件会切成不同块。跨算法仓库\textbf{设计上}复用率是 0\%——像用不同尺子切同一根香肠。
\end{pitfall}
